$\mathbf{B}(q)$ & 第 $q$ 个采样时刻的三轴磁场观测向量 \\
$\bar{\mathbf{B}}(q)$ & 第 $q$ 个采样时刻的无车基线估计向量 \\
$d(q)$ & 第 $q$ 个采样时刻的标量磁场扰动检测统计量 \\
$b_{l,r}$ & 位于边界车道线 $l$ 且属于第 $r$ 组的地磁传感器 \\
$N_{\mathrm{grp}}$ & 沿车辆行驶方向布设的地磁传感器组数 \\
$\mathbf{s}_{l,r}$ & 地磁传感器 $b_{l,r}$ 的空间位置向量 \\
$L_a$ & 相邻地磁传感器组的纵向间距 \\
$\mathbf{z}_j$ & 后台按到达顺序接收的第 $j$ 条量测向量 \\
$t_j$ & 量测 $\mathbf{z}_j$ 的触发时间戳 \\
$b_j$ & 产生量测 $\mathbf{z}_j$ 的地磁传感器编号 \\
$l_j$ & 量测 $\mathbf{z}_j$ 对应的边界车道线编号 \\
$Mpeak_j$ & 量测 $\mathbf{z}_j$ 的地磁扰动峰值 \\
$r(j)$ & 量测 $\mathbf{z}_j$ 所属的观测位置组编号 \\
$\Delta T$ & 扫描时间片长度 \\
$\mathcal{Z}_n$ & 第 $n$ 个扫描内的量测集合 \\
$\mathcal{Z}_n^{\uparrow}$ & 按触发时间升序排序后的第 $n$ 个扫描量测集合 \\
$\mathcal{W}_n$ & 以第 $n$ 个扫描为右端点的固定滞后活动窗口 \\
$\mathcal{S}_{n-L}$ & 固定滞后窗口左端保留的快照状态集合 \\
$\hat{\mathbf{x}}_{i,k|k}$ & 轨迹 $i$ 在时刻 $k$ 的后验状态估计 \\
$\mathbf{P}_{i,k|k}$ & 轨迹 $i$ 在时刻 $k$ 的后验协方差矩阵 \\
$P_{E,i,n}$ & 轨迹 $i$ 在第 $n$ 次扫描时的存在概率 \\
$\boldsymbol{\pi}_{i,n}^{\mathrm{lane}}$ & 轨迹 $i$ 在第 $n$ 次扫描时的车道后验概率向量 \\
$\boldsymbol{\pi}_{r,n}^{\mathrm{scene}}$ & 第 $r$ 组观测位置在第 $n$ 次扫描时的场景后验向量 \\
$\mathbf{x}_{i,k}$ & 轨迹 $i$ 在第 $k$ 次有效更新后的状态向量 \\
$p_{i,k}$ & 轨迹 $i$ 在第 $k$ 次有效更新后的纵向位置 \\
$v_{i,k}$ & 轨迹 $i$ 在第 $k$ 次有效更新后的纵向速度 \\
$\mathbf{F}(\Delta t_{i,k+1})$ & 状态转移矩阵 \\
$\mathbf{Q}(\Delta t_{i,k+1})$ & 过程噪声协方差矩阵 \\
$\tilde{p}_{i,k+1}$ & 轨迹 $i$ 在第 $k+1$ 次更新时的伪位置观测 \\
$\mathbf{H}$ & 观测矩阵 \\
$R_{b_j}$ & 地磁传感器 $b_j$ 对应的观测方差 \\
$\Delta t_{i,k+1}$ & 轨迹 $i$ 相邻两次有效更新之间的时间间隔 \\
$\nu_{i,k+1}$ & 轨迹 $i$ 在第 $k+1$ 次更新时的新息 \\
$S_{i,k+1}$ & 轨迹 $i$ 在第 $k+1$ 次更新时的新息方差 \\
$\mathbf{K}_{i,k+1}$ & 轨迹 $i$ 在第 $k+1$ 次更新时的 Kalman 增益 \\
$\hat{t}_{i\rightarrow b_j}$ & 轨迹 $i$ 预测通过地磁传感器 $b_j$ 的时刻 \\
$\nu^{t}_{i,j}$ & 轨迹 $i$ 与量测 $\mathbf{z}_j$ 的时间残差 \\
$S^{t}_{i,j}$ & 轨迹 $i$ 与量测 $\mathbf{z}_j$ 的时间残差方差 \\
$d^{t\,2}_{i,j}$ & 轨迹 $i$ 与量测 $\mathbf{z}_j$ 的归一化时间距离 \\
$\gamma_t$ & 时间跟踪门阈值 \\
$V^{t}_{i,j}$ & 轨迹 $i$ 对量测 $\mathbf{z}_j$ 的时间跟踪门宽度 \\
$g_{i,j}^{\mathrm{time}}$ & 轨迹 $i$ 与量测 $\mathbf{z}_j$ 的时间一致性似然项 \\
$\phi_{i,j}^{\mathrm{mag}}$ & 轨迹 $i$ 与量测 $\mathbf{z}_j$ 的地磁峰值一致性项 \\
$\mu_i^{M}$ & 轨迹 $i$ 历史对数地磁峰值的均值 \\
$(\sigma_i^{M})^2$ & 轨迹 $i$ 历史对数地磁峰值的方差 \\
$\phi_{i,j}^{\mathrm{lane}}$ & 结合场景后验的车道一致性项 \\
$\Lambda_{i,j}(n)$ & 第 $n$ 次扫描中轨迹 $i$ 与量测 $\mathbf{z}_j$ 的标准化关联似然比 \\
$\beta_{i,j}(n)$ & 第 $n$ 次扫描中轨迹 $i$ 与量测 $\mathbf{z}_j$ 的关联概率 \\
$P_D^{(b)}$ & 地磁传感器 $b$ 的检测概率 \\
$\lambda_{c,b}$ & 地磁传感器 $b$ 的杂波空间密度 \\
$\alpha_{D,b},\beta_{D,b}$ & 地磁传感器 $b$ 检测概率的 Beta 后验参数 \\
$a_{\lambda_c,b},b_{\lambda_c,b}$ & 地磁传感器 $b$ 杂波空间密度的 Gamma 后验参数 \\
$\hat{P}_D^{(b)}(n)$ & 第 $n$ 次扫描后地磁传感器 $b$ 的检测概率后验均值 \\
$\hat{\lambda}_{c,b}(n)$ & 第 $n$ 次扫描后地磁传感器 $b$ 的杂波空间密度后验均值 \\
$\rho_D,\rho_c$ & 检测概率与杂波空间密度更新的遗忘因子 \\
$\mathbf{o}_{r,n}$ & 第 $r$ 组观测位置在第 $n$ 次扫描的场景观测向量 \\
$\Delta t_{r,n}^{\mathrm{adj}}$ & 第 $r$ 组观测位置双侧量测的调整后触发时间差 \\
$\alpha_{r,n}$ & 第 $r$ 组观测位置双侧地磁峰值构成的相对强弱特征 \\
$\mathcal{E}_{r,n}^{(1)}$ & 最近 $L$ 次扫描内单轨迹解释模型的累计代价 \\
$\mathcal{E}_{r,n}^{(2)}$ & 最近 $L$ 次扫描内双轨迹解释模型的累计代价 \\
$\Gamma_{r,n}(s)$ & 场景状态 $s$ 下的轨迹可解释性因子 \\
$\mathcal{I}_1$ & 左车道真实车辆伴随右侧干扰的场景状态 \\
$\mathcal{I}_2$ & 右车道真实车辆伴随左侧干扰的场景状态 \\
$\mathcal{P}$ & 双车并行场景状态 \\
$p_S,p_B$ & 轨迹存活概率与新生概率 \\
$\bar{P}_{D,i}(n)$ & 轨迹 $i$ 在第 $n$ 次扫描预测位置附近的平均检测概率 \\
$N_{\mathrm{miss}}^{(i)}(n)$ & 轨迹 $i$ 在第 $n$ 次扫描的连续漏检计数 \\
$p_c$ & 单步换道概率 \\
$\varepsilon_M$ & 地磁对数变换中的稳定项 \\
$\mathbf{u}_m$ & 第 $m$ 个输入聚类样本 \\
$\mathbf{q}_m$ & 样本 $\mathbf{u}_m$ 的方向向量 \\
$g_m$ & 样本 $\mathbf{u}_m$ 的地磁代表值 \\
$n_m$ & 样本 $\mathbf{u}_m$ 的有效观测数量（支持度） \\
$t_m^{\mathrm{s}},t_m^{\mathrm{e}}$ & 样本 $\mathbf{u}_m$ 的起始与结束时间戳 \\
$\Delta \mathbf{r}_m$ & 样本 $\mathbf{u}_m$ 在位置、车道和时间三维空间中的归一化差分向量 \\
$D_p,D_{\ell},D_t$ & 位置、车道和时间维度的归一化尺度 \\
$\mathbf{q}_{\mathrm{road}}$ & 由道路几何结构和行驶方向确定的单位方向向量 \\
$\varepsilon_q$ & 方向向量归一化中的稳定项 \\
$\mathbb{C}$ & 当前系统中的簇集合 \\
$C_c$ & 第 $c$ 个轨迹重构簇 \\
$\mathrm{Set}_c$ & 簇 $C_c$ 内的节点集合 \\
$\mathrm{Dir}_c$ & 簇 $C_c$ 内的方向约束集合 \\
$\mathcal{A}_c$ & 簇 $C_c$ 的簇尾摘要 \\
$\bar{\mathbf{q}}_c$ & 簇 $C_c$ 当前维护的代表方向向量 \\
$\bar{g}_c$ & 簇 $C_c$ 当前维护的地磁代表值 \\
$N_c$ & 簇 $C_c$ 的累计支持度 \\
$t_c^{\mathrm{tail}}$ & 簇 $C_c$ 尾端样本的结束时间 \\
$E_c$ & 簇 $C_c$ 的活跃度 \\
$s_c$ & 簇 $C_c$ 的生命周期状态 \\
$\Delta t_{cm}$ & 样本 $\mathbf{u}_m$ 起始时间与簇 $C_c$ 尾端时间之间的间隔 \\
$T_{\max}$ & 允许拼接的最大时间间隔 \\
$g_{cm}^{\mathrm{dir}}$ & 样本 $\mathbf{u}_m$ 与簇 $C_c$ 的方向跟踪门函数 \\
$S_{cm}^{\mathrm{dir}}$ & 样本 $\mathbf{u}_m$ 与簇 $C_c$ 的方向相似度 \\
$S_{cm}^{\mathrm{mag}}$ & 样本 $\mathbf{u}_m$ 与簇 $C_c$ 的地磁相似度 \\
$\sigma_g$ & 地磁相似度衰减尺度参数 \\
$\varepsilon_g$ & 地磁相似度计算中的稳定项 \\
$\mathrm{Score}_{cm}$ & 样本 $\mathbf{u}_m$ 与簇 $C_c$ 的聚类相似度评分 \\
$w_q,w_g$ & 方向项与地磁项的评分权重 \\
$\tau_q$ & 方向跟踪门阈值 \\
$\tau_{\mathrm{clus}}$ & 样本归簇接纳阈值 \\
$\tau_E$ & 簇活跃度衰减时间常数 \\
$\alpha_E$ & 簇活跃度增益系数 \\
$n_{\mathrm{ref}}$ & 支持度归一化参考值 \\
$E_{\min}$ & 簇关闭判定的最小活跃度阈值 \\
$T_{\mathrm{del}}$ & 簇关闭或删除的最长未更新时间阈值 \\
$P_{\mathrm{link}}$ & 片段拼接精确率 \\
$R_{\mathrm{link}}$ & 片段拼接召回率 \\
$F1_{\mathrm{link}}$ & 片段拼接 $F1$ 值 \\
$R_{\mathrm{wm}}$ & 错并率 \\
$\mathrm{Frag}_{\mathrm{drop}}$ & 碎片消减量 \\
$\mathrm{FMI}$ & 聚类一致性指标 \\